\documentclass[11pt]{article}
\usepackage[margin=1in]{geometry}
\usepackage{amsmath,amssymb,amsthm,mathtools}
\usepackage{booktabs}
\usepackage{enumitem}
\usepackage{hyperref}
\title{Step 123: Balanced-Length Theorem for the Regularized Müntz/co-Poisson Shadow}
\author{RATCHET / RH Membrane Program}
\date{May 2026}

\newtheorem{theorem}{Theorem}
\newtheorem{lemma}{Lemma}
\newtheorem{definition}{Definition}
\newtheorem{corollary}{Corollary}

\begin{document}
\maketitle

\section{Purpose}
Step 122 made the Burnol-to-Dirichlet readability gate lawful by replacing naive critical-line partial sums with a smooth M\"untz/co-Poisson shadow.  The new problem is quantitative: the M\"untz shadow must be long enough to approximate the Burnol/co-Poisson target on the declared vertical window, but short enough to remain inside the available Heap--Soundararajan style source estimates.

This note gives the first balanced-length theorem.  It does not prove RH.  It states the exact compatibility inequality that must hold before the source route can close.

\section{Objects}
Let
\[
Y_{B,N}
\]
be a finite Burnol/co-Poisson atom window, and let
\[
B_N:Y_{B,N}\to H_N
\]
be the exact Burnol response synthesis map.  Let
\[
D_N^{\rm reg}:\mathbb C^{I_N}\to H_N
\]
be the declared regularized M\"untz/co-Poisson Dirichlet-readable synthesis family, and let
\[
P_N^{\rm reg}
\]
denote orthogonal projection onto its range.  The regularized shadow defect is
\[
\delta_{BD,N}^{\rm reg}
=
\left\|(I-P_N^{\rm reg})B_NG_{B,N}^{-1/2}\right\|.
\]
Equivalently,
\[
\Xi_{BD,N}^{\rm reg}=B_N^*(I-P_N^{\rm reg})B_N.
\]

The effective hybrid visibility constant has the form
\[
c_{{\rm hyb},N}
\ge
1-
\bigl(
\epsilon_{B,N}+\tau_{{\rm reg},N}
\bigr)^2,
\]
where
\[
\tau_{{\rm reg},N}
=
\epsilon_{{\rm Mell},N}+\tau_{M,N}+\kappa_{{\rm tail},N}.
\]

\section{Length parameters}
Let
\[
T_N = \text{vertical / spectral window},
\]
\[
X_N = \text{M\"untz smoothing length},
\]
\[
R_N = \text{Burnol atom intrinsic coefficient length},
\]
\[
A_N = \text{uniform Burnol atom complexity},
\]
\[
Q_N = \text{source conductor / averaging scale}.
\]
After multiplying the Burnol atom shadow by the regularized zeta shadow, the effective Dirichlet length is
\[
L_N\asymp X_NR_N.
\]
The source method is assumed valid only in the short-polynomial range
\[
L_N\le Q_N^{\theta_{\rm src}-o(1)}.
\]
For the direct Heap--Soundararajan template one should think of \(\theta_{\rm src}\) as small, since their construction deliberately keeps the Dirichlet polynomial short.

\section{Error model}
We use the following abstract but checkable model:
\[
\tau_{M,N}\le C A_N(1+T_N)^\alpha X_N^{-\eta},
\]
\[
\epsilon_{{\rm Mell},N}+\kappa_{{\rm tail},N}\le C A_NR_N^{-\beta}.
\]
Here \(\eta>0\) is the decay exponent obtained by moving the M\"untz contour, \(\alpha\) is the vertical-window growth exponent, and \(\beta\) is the Burnol-to-coefficient modeling exponent.

\section{Polynomial exponent form}
Suppose
\[
T_N=Q_N^t,
\quad
A_N\le Q_N^a,
\quad
X_N=Q_N^x,
\quad
R_N=Q_N^r.
\]
Then the conditions become:
\[
a+\alpha t-\eta x<0,
\]
\[
a-\beta r<0,
\]
\[
x+r<\theta_{\rm src}.
\]

\begin{theorem}[Balanced-length compatibility]
Under the polynomial error model above, there exist exponents \(x,r>0\) such that the regularized shadow errors vanish and the Dirichlet polynomial remains in the source-valid range if and only if
\[
\boxed{
\frac{a+\alpha t}{\eta}+\frac{a}{\beta}<\theta_{\rm src}.
}
\]
More precisely, any choice satisfying
\[
x>\frac{a+\alpha t}{\eta},
\qquad
r>\frac{a}{\beta},
\qquad
x+r<\theta_{\rm src}
\]
works, and no polynomial choice works if the displayed strict inequality fails.
\end{theorem}

\begin{proof}
The M\"untz contour residual tends to zero if
\[
A_NT_N^\alpha X_N^{-\eta}=Q_N^{a+\alpha t-\eta x}\to0,
\]
which is equivalent to \(x>(a+\alpha t)/\eta\).  The Mellin/coefficient and support-tail errors tend to zero if
\[
A_NR_N^{-\beta}=Q_N^{a-\beta r}\to0,
\]
which is equivalent to \(r>a/\beta\).  The source estimate remains lawful if
\[
X_NR_N=Q_N^{x+r}\le Q_N^{\theta_{\rm src}-o(1)},
\]
which requires \(x+r<\theta_{\rm src}\).  These three strict inequalities are feasible exactly when the sum of the two lower bounds is less than \(\theta_{\rm src}\).
\end{proof}

\section{Consequences}
\subsection{Same-aspect obstruction}
If \(Q_N\) and \(T_N\) are the same scale, then \(t=1\).  With bounded atom complexity \(a=0\) and no intrinsic coefficient length cost \(r=0\), the necessary condition is
\[
\frac{\alpha}{\eta}<\theta_{\rm src}.
\]
This is the simplest obstruction.  If the contour error grows like a positive power of the same vertical window and the source machinery only allows very short polynomials, the same-aspect route is blocked.

\subsection{Decoupled conductor route}
If the source conductor scale \(Q_N\) may grow faster than the vertical window \(T_N\), then \(t<1\).  The condition becomes
\[
\frac{a+\alpha t}{\eta}+\frac{a}{\beta}<\theta_{\rm src}.
\]
Thus a sufficiently large conductor family can leave room for both M\"untz accuracy and source validity.  This points toward the conductor/aspect character ladder rather than a pure t-aspect route.

\subsection{Positive floor regime}
Vanishing readability error is sufficient but not strictly necessary for source absorption.  If
\[
\epsilon_{B,N}+\tau_{{\rm reg},N}\le 1-\delta
\]
for some \(\delta>0\), then
\[
c_{{\rm hyb},N}\ge 1-(1-\delta)^2>0.
\]
Hence any source strength \(\gamma_N\to\infty\) gives
\[
\gamma_Nc_{{\rm hyb},N}\to\infty.
\]
This still requires fixed/exhaustive tail promotion; a positive visibility floor alone is not an RH proof.

\section{RH-route reading}
The source route now has five quantitative obligations:
\[
\epsilon_{B,N}\quad\text{Burnol geometric exhaustion},
\]
\[
\epsilon_{{\rm Mell},N}\quad\text{Mellin/Dirichlet modeling},
\]
\[
\tau_{M,N}\quad\text{M\"untz contour residual},
\]
\[
\kappa_{{\rm tail},N}\quad\text{finite-window tail/support},
\]
\[
\gamma_N\quad\text{Hecke/Dirichlet source strength}.
\]
The accepted lower-frame route requires
\[
\boxed{
\gamma_N\left[1-(\epsilon_{B,N}+\epsilon_{{\rm Mell},N}+\tau_{M,N}+\kappa_{{\rm tail},N})^2\right]\to\infty
}
\]
with fixed/exhaustive ledger promotion.

\section{Strategic verdict}
The same-aspect version is likely underpowered if one imports the current short-polynomial moment technology without improvement.  The viable route is therefore:
\[
\boxed{
\text{decouple the source conductor scale from the vertical M\"untz window.}
}
\]
Equivalently, use a conductor/aspect Hecke or Dirichlet family to obtain a larger source-valid range than the M\"untz shadow length would permit in pure t-aspect.

\section{Next target}
The next concrete theorem is a conductor-decoupled source theorem:
\[
\boxed{
\text{construct }Q_N\gg T_N\text{ source families such that }L_N=X_NR_N\le Q_N^{\theta_{\rm src}}
}
\]
while retaining the residual-specific lower-frame inequality on \(\Xi^{\rm BC}\).

\end{document}
